% =====================================================================
% CODEX ColombiaMacro · Parte IV · Interpretar en conjunto (español)
% Fragmento para \input: sin preámbulo.
% =====================================================================
\providecommand{\cmsym}[1]{{\fontspec{DejaVu Sans}#1}}

\part{Interpretar en conjunto}

% ---------------------------------------------------------------------
\chapter{Cómo se conectan las piezas}
\label{es:i:cadena}

Un gráfico aislado responde una pregunta; una economía se entiende cuando las respuestas se encadenan. Esta parte
del Codex explica cómo leer las once páginas como un solo sistema. No describe lo que \emph{va} a pasar: describe
qué relaciones conceptuales unen a los indicadores y cómo comprobar, con los datos del sitio, si esas relaciones
se están cumpliendo o no en un momento dado.

\section{La cadena básica}

La secuencia de lectura del sitio sigue el orden en que los choques se transmiten en una economía pequeña y
abierta con un banco central que fija una meta de inflación:

\begin{enumerate}
\item \textbf{Actividad y capacidad} (Ciclo, Crecimiento, Capacidad). Si la economía produce por encima de lo que
  puede producir de manera sostenida (brecha positiva), las empresas usan más intensamente su capital y su personal.
\item \textbf{Mercado laboral} (Empleo). Una brecha positiva suele acompañarse de desempleo por debajo de su
  tendencia (ley de Okun) y de presión sobre los salarios, sobre todo en los servicios.
\item \textbf{Precios} (Inflación). La presión de demanda aparece primero en la inflación de servicios y en la
  inflación de fondo; los alimentos y los regulados responden más a choques de oferta, al clima y a decisiones
  administrativas.
\item \textbf{Reacción de política} (Tasas). El Banco de la República ajusta su tasa para que la tasa \emph{real}
  (descontada la inflación esperada) quede por encima de la neutral cuando quiere enfriar la demanda, o por debajo
  cuando quiere estimularla.
\item \textbf{Precios de los activos} (Curva TES, Tasa de cambio, Empresas). El mercado incorpora esa reacción en
  la curva de rendimientos (la parte corta sigue a la tasa de política; la larga mezcla expectativas y primas de
  riesgo), en el peso (diferencial de tasas y apetito por riesgo) y en la bolsa (costo del capital y utilidades).
\item \textbf{Restricciones} (Externo y fiscal, Comercio). El país financia la diferencia entre lo que gasta y lo que
  produce con ahorro externo (déficit de cuenta corriente) y el Gobierno financia su déficit con deuda. Estas
  restricciones determinan cuánto margen tienen la política monetaria y la fiscal, y cuán sensible es el país a
  cambios en el financiamiento internacional y en los precios de sus exportaciones.
\end{enumerate}

\begin{leer}[La cadena no es una ley]
Cada eslabón es una relación \emph{típica}, no automática. Los choques de oferta (clima, precios internacionales,
regulación) pueden mover la inflación sin que cambie la brecha; un choque externo puede mover el peso sin que
cambie la política monetaria. Por eso el sitio muestra cada eslabón por separado: para que el lector compruebe,
dato por dato, si la cadena se está cumpliendo.
\end{leer}

\section{Mapa de relaciones entre páginas}

\begin{center}
\small
\begin{tabularx}{\linewidth}{>{\raggedright\arraybackslash}p{3.3cm}>{\raggedright\arraybackslash}X>{\raggedright\arraybackslash}p{3.6cm}}
\toprule
\textbf{Si observa\ldots} & \textbf{\ldots compruebe en} & \textbf{Para saber si\ldots}\\
\midrule
Brecha del producto positiva & Empleo (desempleo frente a su tendencia, subutilización), Inflación (servicios y fondo) &
la presión de demanda llega al mercado laboral y a los precios\\
Inflación total sube & Inflación (aportes por división, regulados y alimentos, difusión) & el alza es amplia o se
concentra en pocos rubros\\
Tasa de política sube & Tasas (traspaso a IBR, CDT y crédito; crédito real) & la decisión se transmite al sistema
financiero\\
Pendiente de la curva baja o se invierte & Curva TES (descomposición tasa real e inflación implícita; prima sobre
la TPM) & el movimiento viene de expectativas de inflación o de tasa real\\
Peso se deprecia & Tasa de cambio (pares regionales, dólar global, Brent, volatilidad) & el movimiento es propio de
Colombia o global\\
Bolsa sube & Empresas (equiponderado frente a COLCAP, 7 Magníficas, amplitud) & el alza es amplia o concentrada en
pocas empresas\\
Déficit de cuenta corriente se amplía & Externo (componentes, IED, financiación) y Comercio (precio frente a
volumen) & el déficit está financiado con flujos estables y si responde a precios o a cantidades\\
Déficit fiscal se amplía & Externo y fiscal (ingresos, gastos, intereses, balance primario, deuda) & el deterioro
viene de menores ingresos, mayor gasto o mayor costo de la deuda\\
\bottomrule
\end{tabularx}
\end{center}

\section{Tres velocidades de la información}

Los indicadores del sitio llegan con rezagos distintos, y la lectura conjunta debe respetarlo:

\begin{itemize}
\item \textbf{Diarios} (TRM, COLCAP, curva TES, tasa de política, IBR): reaccionan en el mismo día a la información
  nueva; son ruidosos y deben leerse con promedios o cambios de varios meses.
\item \textbf{Mensuales} (IPC, GEIH, ISE, comercio exterior, crédito, remesas): muestran el estado de la economía con
  uno o dos meses de rezago; tienen estacionalidad y se leen en variación anual o en sumas de 12 meses.
\item \textbf{Trimestrales y anuales} (PIB, cuenta corriente, cuentas departamentales, 10.000 empresas, deuda del
  Gobierno): son las mediciones más completas pero llegan con uno a dieciocho meses de rezago y se revisan.
\end{itemize}

\begin{alerta}
No compare directamente un dato diario de hoy con un dato trimestral de hace seis meses como si describieran el mismo
momento. Cada cifra del sitio lleva su fecha: úsela para ubicar cada pieza en el tiempo.
\end{alerta}

% ---------------------------------------------------------------------
\chapter{Rutina de lectura mensual}
\label{es:i:rutina}

La siguiente rutina, de unos treinta minutos, recorre el sitio en el orden de la cadena y deja registradas las
respuestas clave. Está pensada para una revisión mensual, después de que el DANE publica el IPC y la GEIH.

\section{Paso 1 · Portada (cinco minutos)}
\begin{enumerate}
\item Lea el recuadro del ciclo: fase, brecha, dirección y racha.
\item Recorra las doce tarjetas de ``Lectura en un vistazo'' y anote los estados que cambiaron frente al mes
  anterior (por ejemplo, de ``En meta'' a ``Sobre la meta'').
\item En ``Todos los indicadores'', identifique los que están en percentiles extremos (por debajo de 10 o por
  encima de 90): son los datos inusuales frente a su propia historia.
\end{enumerate}

\section{Paso 2 · Actividad (cinco minutos)}
\begin{enumerate}
\item \textbf{Ciclo:} ¿cuántos de los cinco métodos ven la economía por encima de su capacidad? ¿Cuántos de los doce
  sectores están sobre su tendencia (amplitud)?
\item \textbf{Crecimiento:} compare las tres velocidades (anual, 12 meses, trimestral anualizada). Si la trimestral
  anualizada está por debajo de la anual, el ritmo reciente es menor que el del último año.
\item \textbf{Capacidad:} revise si la brecha laboral (desempleo frente a su tendencia) confirma la brecha del
  producto.
\end{enumerate}

\section{Paso 3 · Empleo e inflación (diez minutos)}
\begin{enumerate}
\item \textbf{Empleo:} desempleo desestacionalizado, ocupados frente a un año antes, informalidad y por qué ramas se
  crea el empleo.
\item \textbf{Inflación:} total frente al rango meta; de fondo (sin alimentos ni regulados); servicios; difusión
  (proporción de subclases por encima de 4\%); aportes por división.
\item Pregunta clave: \emph{¿la inflación de fondo y la de servicios se mueven en la misma dirección que la
  brecha?} Si sí, la presión es de demanda; si la total se mueve sola, el origen está en alimentos, regulados o
  energía.
\end{enumerate}

\section{Paso 4 · Política y mercados (cinco minutos)}
\begin{enumerate}
\item \textbf{Tasas:} tasa real ex ante frente al rango neutral; traspaso a las tasas del sistema financiero;
  crecimiento real del crédito.
\item \textbf{Curva TES:} nivel, pendiente y forma; inflación implícita a 10 años y la 5y5y frente al techo del
  rango meta (4\%).
\item \textbf{Tasa de cambio:} variación de 12 meses del peso frente a sus pares; dólar global; Brent; volatilidad
  frente a su percentil 80.
\item \textbf{Empresas:} COLCAP frente al equiponderado; utilidades y endeudamiento de las 10.000 empresas; crédito
  a empresas.
\end{enumerate}

\section{Paso 5 · Restricciones (cinco minutos)}
\begin{enumerate}
\item \textbf{Externo:} déficit de cuenta corriente de cuatro trimestres y qué parte cubre la IED; remesas; deuda
  externa y posición de inversión internacional.
\item \textbf{Fiscal:} balance total y primario del Gobierno en 12 meses; intereses sobre ingresos; deuda bruta.
\item \textbf{Comercio:} exportaciones e importaciones de 12 meses; cuánto del cambio en el valor es precio y cuánto
  volumen; concentración por destinos y productos.
\end{enumerate}

\begin{dato}[Registro sugerido]
Al final de la rutina, escriba cinco frases, una por eslabón de la cadena (actividad, precios, política, mercados,
restricciones), cada una con la cifra y la fecha que la respaldan. Es exactamente la estructura del veredicto
automático del sitio, y permite comparar su propia lectura con la regla explícita.
\end{dato}

% ---------------------------------------------------------------------
\chapter{Lista de preguntas para el inversionista}
\label{es:i:preguntas}

Las preguntas siguientes se pueden responder \emph{solo} con información del sitio. Están agrupadas por el tipo de
decisión que suelen informar. Responderlas no equivale a una recomendación: son el insumo para un análisis propio.

\section{Riesgo macroeconómico general}
\begin{itemize}
\item ¿La economía está por encima o por debajo de su capacidad, y cuán robusta es esa lectura entre métodos?
\item ¿El crecimiento es amplio (muchos sectores y regiones) o concentrado (pocos sectores, sector público)?
\item ¿El desempleo está por debajo o por encima de su tendencia? ¿La informalidad baja o sube?
\end{itemize}

\section{Riesgo de tasas e inflación}
\begin{itemize}
\item ¿La inflación está dentro del rango meta de 2\% a 4\%? ¿Acelera, cede o está estable en los últimos tres meses?
\item ¿Cuántas subclases del IPC están por encima de 4\%? ¿Qué divisiones explican la mayor parte de la inflación?
\item ¿La tasa real ex ante está por encima o por debajo del rango neutral? ¿Cuánto se ha transmitido el último ciclo
  de la tasa de política a las tasas de crédito?
\item ¿Las expectativas de largo plazo implícitas en los TES (5y5y) están por debajo del techo de 4\%?
\item ¿La curva tiene pendiente normal, plana o invertida? ¿Qué parte del movimiento reciente es tasa real y qué
  parte inflación implícita?
\end{itemize}

\section{Riesgo cambiario}
\begin{itemize}
\item ¿El movimiento del peso en 12 meses es propio de Colombia o lo comparten el real, el peso mexicano y el sol?
\item ¿El peso está caro o barato en términos reales frente a su historia (ITCR)?
\item ¿Cómo se comporta el dólar global y el Brent, dos determinantes externos del peso?
\item ¿Qué tan volátil está el peso frente a su propia historia?
\end{itemize}

\section{Riesgo externo y fiscal}
\begin{itemize}
\item ¿El déficit de cuenta corriente supera el 4\% del PIB? ¿Qué proporción cubre la inversión extranjera directa,
  el flujo más estable?
\item ¿La deuda externa y la posición de inversión internacional neta se deterioran o mejoran como proporción del PIB?
\item ¿El balance primario del Gobierno es positivo o negativo? ¿Qué proporción de los ingresos se va en intereses?
\item ¿La deuda bruta del Gobierno está por encima de 60\% del PIB, el nivel de referencia que marca el sitio?
\end{itemize}

\section{Riesgo sectorial y empresarial}
\begin{itemize}
\item ¿Qué sectores crecen en producción y cuáles en empleo? ¿Coinciden?
\item ¿El alza de la bolsa es amplia (equiponderado acompaña al COLCAP) o concentrada en las 7 Magníficas?
\item ¿Las utilidades, el margen y la rentabilidad del patrimonio de las 10.000 empresas mejoran o se deterioran?
  ¿Sube el endeudamiento?
\item ¿Crece en términos reales el crédito a las empresas? ¿Se crean más sociedades de las que se cancelan?
\item ¿El valor exportado sube por precio o por volumen? ¿Cuán concentradas están las exportaciones por destino y
  producto?
\end{itemize}

% ---------------------------------------------------------------------
\chapter{Los estados automáticos y el veredicto}
\label{es:i:estados}

El sitio no usa juicios subjetivos para etiquetar los datos. Cada etiqueta (``Alta'', ``Restrictiva'', ``Peso
fuerte'') sale de una regla fija sobre el último dato. Conocer las reglas permite saber exactamente qué significa
cada etiqueta y cuándo cambiará.

\section{Parámetros}
\begin{center}
\small
\begin{tabularx}{\linewidth}{lX}
\toprule
\textbf{Parámetro} & \textbf{Valor y significado}\\
\midrule
Rango meta de inflación & $3\% \pm 1$ pp (Banco de la República)\\
Tasa real neutral & entre $2{,}7\%$ y $3{,}0\%$ real (estimación externa del equipo técnico del Banco)\\
Margen de postura & $0{,}5$ pp alrededor del rango neutral\\
Umbral de brecha & $\lvert g\rvert<0{,}5\%$: ``cerca del potencial''\\
Ancla de expectativas & la 5y5y por encima de $4\%$ indica expectativas desancladas\\
\bottomrule
\end{tabularx}
\end{center}

\section{Reglas del veredicto técnico}

Sea $g$ la brecha del producto (en \% del potencial), $\pi_t$ la inflación anual, $r^{ea}$ la tasa real ex ante y
$\pi^{5y5y}$ la inflación implícita entre 5 y 10 años:

\begin{center}
\small
\begin{tabularx}{\linewidth}{>{\raggedright\arraybackslash}p{3.4cm}X>{\raggedright\arraybackslash}p{4cm}}
\toprule
\textbf{Concepto} & \textbf{Regla} & \textbf{Resultado}\\
\midrule
Posición frente al potencial & $g\ge 0{,}5$; $g\le-0{,}5$; en otro caso & por encima; por debajo; cerca\\
Fase del ciclo & signo de $g$ y de $\Delta_2 g = g_t-g_{t-2}$ & expansión, desaceleración, contracción, recuperación\\
Dirección de la inflación & $\pi_t-\pi_{t-3}>0{,}1$; $<-0{,}1$; en otro caso & acelerando; cediendo; estable\\
Rango meta & $\pi_t>4$ o $\pi_t<2$ & fuera del rango; si no, dentro\\
Postura monetaria & $r^{ea}>3{,}5$; $r^{ea}<2{,}2$; en otro caso & restrictiva; expansiva; neutral\\
Expectativas & $\pi^{5y5y}>4$ & desancladas; si no, ancladas\\
\bottomrule
\end{tabularx}
\end{center}

\section{Estados de la portada}

\begin{center}
\small
\begin{tabularx}{\linewidth}{>{\raggedright\arraybackslash}p{2.6cm}X>{\raggedright\arraybackslash}p{4.2cm}}
\toprule
\textbf{Tarjeta} & \textbf{Regla} & \textbf{Estados}\\
\midrule
Crecimiento & PIB anual $g^Y$ frente al crecimiento potencial $g^*$: $g^Y<0$; $g^Y\ge g^*+1$; $g^Y\ge g^*-1$; en otro caso & contracción; fuerte (por encima de lo habitual); normal; lento\\
Inflación & $2\le\pi\le4$; $\pi<2$; $\pi>5$; $4<\pi\le5$ & en meta; baja; alta; sobre la meta\\
Desempleo & frente a 12 meses antes: baja más de $0{,}3$ pp; sube más de $0{,}3$ pp; en otro caso & mejorando; empeorando; estable\\
Tasa del Banco & postura del veredicto & restrictiva (frena la economía), neutral, expansiva\\
Dólar & variación de 12 meses de la TRM: $<-5\%$; $>5\%$; en otro caso & peso fuerte; peso débil; estable\\
Bolsa & variación de 12 meses del COLCAP: $>5\%$; $<-5\%$; en otro caso & sube; baja; estable\\
Nivel histórico & percentil $<10$; $<30$; $\le70$; $\le90$; $>90$ & muy bajo; bajo; normal; alto; muy alto\\
\bottomrule
\end{tabularx}
\end{center}

\begin{leer}[Qué significan y qué no]
Un estado describe la posición del último dato frente a una regla; no es una calificación de la política económica
ni una señal de compra o venta. ``Alta: frena la economía'' significa que la tasa real supera el rango neutral más el
margen, no que la tasa ``deba'' bajar. ``Peso fuerte'' significa que el dólar cuesta al menos 5\% menos que hace un
año, sin juicio sobre si eso es bueno o malo.
\end{leer}

\section{Percentiles: ``normal'' frente a su propia historia}

La columna de nivel histórico de ``Todos los indicadores'' ubica el último dato en la distribución de la serie desde
2010:
\[
P = 100\times\frac{\#\{s : x_s \le x_t\}}{N}.
\]
Un percentil de 92 significa que el dato actual es mayor o igual que el 92\% de los datos desde 2010. Es una medida
descriptiva: no dice si el nivel es sostenible ni hacia dónde irá. Un indicador puede pasar años en ``muy alto'' si
la economía cambió de régimen.

% ---------------------------------------------------------------------
\chapter{Combinaciones de señales}
\label{es:i:combinaciones}

Las combinaciones siguientes describen situaciones \emph{conceptuales} que la literatura macroeconómica distingue.
Sirven para organizar la lectura, no para anticipar lo que ocurrirá: en cada caso se indica qué gráficos permiten
comprobar cada elemento.

\section{Sobrecalentamiento}
\textbf{Señales:} brecha positiva y amplia (mayoría de métodos y sectores), desempleo por debajo de su tendencia,
inflación de fondo y de servicios por encima de la meta y en aumento, difusión alta.\\
\textbf{Lectura conceptual:} la demanda supera la capacidad. En este escenario la literatura de metas de inflación
describe una tasa real por encima de la neutral como la respuesta típica de política.\\
\textbf{Dónde comprobarlo:} Ciclo (consenso y amplitud), Capacidad (brecha laboral), Inflación (fondo, servicios,
difusión), Tasas (tasa real).

\section{Choque de oferta}
\textbf{Señales:} inflación total sube mientras la de fondo se mantiene; aportes concentrados en alimentos, regulados o
energía; brecha cercana a cero o negativa.\\
\textbf{Lectura conceptual:} los precios suben por costos y no por exceso de demanda. Los efectos directos de un choque
de oferta suelen ser transitorios; el riesgo relevante para un banco central es que contagien a las expectativas
(vigile la 5y5y y la inflación de servicios).\\
\textbf{Dónde comprobarlo:} Inflación (aportes por división, regulados y alimentos), Curva TES (inflación implícita).

\section{Desaceleración con inflación persistente}
\textbf{Señales:} brecha que cae (fase de desaceleración), crecimiento trimestral anualizado por debajo del anual,
inflación todavía fuera del rango meta, tasa real alta.\\
\textbf{Lectura conceptual:} la política monetaria restrictiva actúa con rezagos largos y variables sobre la actividad
y más tarde sobre la inflación. El sitio permite seguir ambos rezagos con la tasa real y con la inflación de fondo.\\
\textbf{Dónde comprobarlo:} Crecimiento (tres velocidades), Inflación (fondo), Tasas (ciclos y traspaso).

\section{Choque externo}
\textbf{Señales:} el peso se deprecia más que sus pares, sube la volatilidad, la curva TES se empina en la parte larga,
cae la bolsa en dólares; puede coincidir con caída del Brent o con un dólar global fuerte.\\
\textbf{Lectura conceptual:} un endurecimiento de las condiciones financieras internacionales o una caída de los
términos de intercambio reduce el ingreso del país y encarece su financiamiento. Si el movimiento es común a la región,
el origen es global; si es propio de Colombia, el origen es local.\\
\textbf{Dónde comprobarlo:} Tasa de cambio (peso frente a vecinos, dólar global, Brent, volatilidad), Curva TES
(prima a 10 años), Externo (financiación de la cuenta corriente).

\section{Bonanza de términos de intercambio}
\textbf{Señales:} precios de exportación en dólares suben más que los de importación, el valor exportado crece por
precio y no por volumen, el peso se aprecia, mejora la cuenta corriente.\\
\textbf{Lectura conceptual:} el ingreso nacional aumenta sin que aumente necesariamente la producción. La literatura
sobre ``monedas de materias primas'' (Chen y Rogoff, 2003) documenta la relación entre precios de exportación y tasa de
cambio real.\\
\textbf{Dónde comprobarlo:} Comercio (precio frente a volumen, términos de intercambio), Tasa de cambio (ITCR, Brent).

\section{Estrés fiscal}
\textbf{Señales:} balance primario negativo y creciente, intereses que absorben una proporción creciente de los
ingresos, deuda bruta por encima de 60\% del PIB, curva TES con mayor prima en la parte larga.\\
\textbf{Lectura conceptual:} la dinámica de la deuda depende de la diferencia entre la tasa de interés real y el
crecimiento ($r-g$) y del balance primario:
\[
\Delta d_t \approx \frac{r_t-g_t}{1+g_t}\,d_{t-1} - pb_t,
\]
donde $d$ es la deuda sobre PIB y $pb$ el balance primario sobre PIB. Con $r>g$, estabilizar la deuda requiere un
superávit primario.\\
\textbf{Dónde comprobarlo:} Externo y fiscal (balance, intereses, deuda), Curva TES (pendiente y prima).

\begin{alerta}
Ninguna combinación ``predice''. Las mismas señales han precedido desenlaces distintos según el contexto. El valor de
este capítulo es ordenar las preguntas y señalar qué datos hay que mirar juntos, no anticipar resultados.
\end{alerta}

% ---------------------------------------------------------------------
\chapter{Errores frecuentes de interpretación}
\label{es:i:errores}

\section{Confundir porcentaje con puntos porcentuales}
Si la inflación pasa de 5\% a 6\%, subió \emph{1 punto porcentual} (pp), no 1\%. En términos relativos subió 20\%:
$(6-5)/5=0{,}20$. El sitio usa ``pp'' para cambios de tasas y ``\%'' para cambios de niveles.

\section{Leer un solo mes}
Una variación mensual alta puede deberse a estacionalidad, a un ajuste administrativo o a un error de muestreo. Las
lecturas del sitio se apoyan en variaciones anuales, sumas de 12 meses o promedios de tres meses por esa razón.

\section{Efecto base}
La variación anual compara con el mismo mes del año anterior. Si hace un año hubo un salto puntual, la variación anual
cae cuando ese salto sale de la comparación, aunque el nivel de este mes no haya cambiado:
\[
\pi_t \approx \sum_{k=0}^{11} m_{t-k},
\]
con $m$ la variación mensual: al entrar un mes nuevo sale el de hace doce meses.

\section{Ilusión nominal}
Un crédito que crece 10\% nominal con inflación de 9\% crece apenas $\approx 0{,}9\%$ real:
$(1{,}10/1{,}09)-1$. Utilidades, salarios, crédito y gasto público deben leerse en términos reales para saber si hay
más bienes y servicios detrás.

\section{Correlación no es causalidad}
Que dos series se muevan juntas (el peso y el Brent, por ejemplo) no prueba que una cause a la otra. Ambas pueden
responder a un tercer factor, como el apetito global por riesgo. El sitio muestra relaciones descriptivas y no estima
efectos causales.

\section{Olvidar las revisiones}
El PIB, el ISE y las cuentas departamentales se revisan. Un dato del último trimestre puede cambiar varias décimas en
la siguiente publicación. Conclusiones que dependen de una décima son frágiles.

\section{Sesgo de supervivencia}
Aplicar la canasta actual del COLCAP hacia atrás, o comparar las 10.000 empresas más grandes de dos años, excluye a las
empresas que salieron de la lista. Esos ejercicios describen a las sobrevivientes, no a todo el universo.

\section{Mezclar precio y retorno total}
El COLCAP es un índice de precios: no incluye dividendos. El retorno de un inversionista que reinvierte dividendos es
mayor que la variación del índice.

\section{Tomar la inflación implícita como pronóstico}
La inflación implícita en los TES incluye primas por riesgo inflacionario y por liquidez. Es ``lo que descuenta el
mercado'', no una expectativa pura ni un pronóstico.

\section{Comparar razones al PIB de fuentes distintas}
La cuenta corriente o la deuda como \% del PIB pueden diferir entre fuentes según el PIB de referencia y la tasa de
cambio usada para convertirlo. Compare razones calculadas de la misma forma.

% ---------------------------------------------------------------------
\chapter{Lo que el sitio no hace}
\label{es:i:limites}

\begin{itemize}
\item \textbf{No proyecta ni pronostica.} Todas las cifras son observadas o se calculan con precios observados. Lo que
  mira hacia adelante (inflación implícita, forward 5y5y) es una lectura del mercado y se identifica como tal.
\item \textbf{No recomienda inversiones.} Los estados y el veredicto describen datos frente a reglas explícitas; no son
  señales de compra o venta ni asesoría financiera.
\item \textbf{No sustituye las publicaciones oficiales.} Ante cualquier diferencia, prevalece la cifra publicada por
  la entidad fuente (DANE, Banco de la República, Ministerio de Hacienda, Supersociedades, Confecámaras).
\item \textbf{No estima efectos causales.} Las relaciones mostradas son descriptivas.
\end{itemize}

\begin{dato}[Uso académico]
Para citar una cifra del sitio en un trabajo académico, cite la fuente oficial que aparece al pie del gráfico, la fecha
del dato y la fecha de consulta. La metodología de cada transformación está en la Nota metodológica.
\end{dato}
